\chapter{Milestone M1.1: Local CAR Algebra, Representations, and Physical Projection}
\label{ch:m1_1}

% EDITORIAL QUESTION: Why must the local spin-1/2 degree of freedom be fractionalized into four Majorana fermions, and what redundancy does this introduce?
\section{Physical Motivation and Spin-1/2 Representation Theory}
\label{sec:m1_1_motivation}
In quantum mechanics, a spin-$S$ degree of freedom is fundamentally defined as an abstract representation of the Lie algebra $\mathfrak{su}(2)$ on a complex Hilbert space $\mathcal{H}$. For spin $S=1/2$, the operator algebra is generated by three Hermitian operators $\mathbf{S} = (S^x, S^y, S^z)$ satisfying the canonical $\mathfrak{su}(2)$ commutation relations:
\begin{equation}
\label{eq:m1_1_spin_su2}
[S^\alpha, S^\beta] = \ii \epsilon^{\alpha\beta\gamma} S^\gamma, \qquad \alpha, \beta, \gamma \in \{x, y, z\},
\end{equation}
subject to the irreducible $\mathrm{SU}(2)$ quadratic Casimir constraint:
\begin{equation}
\label{eq:m1_1_spin_casimir}
\mathbf{S}^2 = (S^x)^2 + (S^y)^2 + (S^z)^2 = S(S+1)\Id = \frac{3}{4}\Id.
\end{equation}
The standard, faithful representation is given on $\mathcal{H}_{\mathrm{phys}} = \mathbb{C}^2$ by the Pauli spin matrices $S^\gamma = \frac{1}{2}\sigma^\gamma$.

In strongly correlated and geometrically frustrated quantum magnets—such as the Kitaev honeycomb model—bond-directional exchange interactions lead to macroscopic classical degeneracies, absence of long-range magnetic ordering down to zero temperature, and topological quantum spin liquid (QSL) states. In such phases, physical spin flips fractionalize into emergent deconfined partons and gauge fluxes.

However, the standard Pauli representation on $\mathbb{C}^2$ does not permit a natural factorization of the spin operator into independent emergent quasiparticles. To access fractionalized phases analytically and variationally, one introduces **parton representations** (such as Abrikosov pseudo-fermions or Schwinger bosons), which represent spin operators as bilinears of canonical creation and annihilation operators. The inevitable price of any parton construction is that it acts on an enlarged Fock space $\mathcal{F}$ with dimension $\dim \mathcal{F} > 2$, introducing an unphysical subspace and an exact local gauge redundancy.

In this chapter, we develop the exact algebraic foundations of the four-Majorana representation per site introduced by Kitaev. To establish mathematical clarity and avoid clutter, **we omit the lattice site index $i$ in Sections~\ref{sec:m1_1_car_algebra} through \ref{sec:m1_1_pauli_algebra}**, since all Clifford, Casimir, and parity identities are strictly on-site (local). We reinstate lattice and bond indices in Section~\ref{sec:m1_1_exchange_tensor} when constructing the two-site exchange bond and the lattice Hamiltonian.

% EDITORIAL QUESTION: How is the local Majorana CAR algebra represented in terms of canonical Dirac fermions on the Fock space?
\section{The Self-Dual CAR Algebra and Local Fock Space}
\label{sec:m1_1_car_algebra}
% SOURCE [VERIFIED]: Kitaev (2006), Ann. Phys. 321, 2-111, Sec. 2 — Self-dual CAR algebra
% CLAIM: K-001 (provenance: doc/rnd/phase_1/m1_1_t1_local_majorana_algebra_conventions.md, test: test/unit/test_local_majorana_algebra.py::test_t00a_01_majorana_car)

\begin{definition}[Self-Dual Canonical Anticommutation Relations]
\label{def:m1_1_car_algebra}
Let $V \cong \mathbb{R}^4$ be a 4-dimensional Euclidean vector space equipped with the standard positive-definite inner product. The self-dual CAR algebra $\mathcal{A}_{\mathrm{CAR}}(V)$ is the real $C^*$-algebra generated by four Hermitian Majorana operators $c = (c^0, c^1, c^2, c^3)^T$ satisfying:
\begin{equation}
\label{eq:m1_1_car_anticommutator}
\{c^\mu, c^\nu\} \equiv c^\mu c^\nu + c^\nu c^\mu = 2 \delta^{\mu\nu} \Id, \qquad (c^\mu)^\dagger = c^\mu, \qquad \forall \mu, \nu \in \{0, 1, 2, 3\}.
\end{equation}
In particular, every Majorana generator squares to the identity: $(c^\mu)^2 = \Id$.
\end{definition}

To construct the explicit Hilbert space representation, we group the four Majorana operators into two canonical complex Dirac fermions:
\begin{equation}
\label{eq:m1_1_dirac_fermions}
d_1 = \frac{1}{2}(c^0 + \ii c^3), \qquad d_2 = \frac{\ii}{2}(c^1 + \ii c^2).
\end{equation}
By direct evaluation using Eq.~\eqref{eq:m1_1_car_anticommutator}, $d_a$ satisfy the canonical Fermi-Dirac algebra:
\begin{equation}
\{d_a, d_b^\dagger\} = \delta_{ab}\Id, \qquad \{d_a, d_b\} = 0, \qquad \forall a, b \in \{1, 2\}.
\end{equation}
The local unprojected Fock space $\mathcal{F} \cong \mathbb{C}^4$ is spanned by the occupation basis:
\begin{equation}
\label{eq:m1_1_fock_basis}
\mathcal{F} = \operatorname{span}\left\{ |00\rangle, \, |01\rangle, \, |10\rangle, \, |11\rangle \right\}, \qquad |n_1, n_2\rangle = (d_1^\dagger)^{n_1} (d_2^\dagger)^{n_2} |00\rangle.
\end{equation}
The dimension of the local Fock space is $\dim \mathcal{F} = 2^{4/2} = 4 = 2 \times \dim \mathcal{H}_{\mathrm{phys}}$, exactly twice the physical spin-$1/2$ dimension.

\begin{theorem}[Four-Majorana Algebraic Contraction Identity]
\label{thm:m1_1_four_majorana_contraction}
Let $c^0, c^1, c^2, c^3$ be four generators of $\mathcal{A}_{\mathrm{CAR}}(V)$. The product of any four Majorana operators on the site expands identically as:
\begin{equation}
\label{eq:m1_1_four_majorana_contraction}
c^\mu c^\nu c^\rho c^\lambda = (\delta^{\mu\nu}\delta^{\rho\lambda} - \delta^{\mu\rho}\delta^{\nu\lambda} + \delta^{\mu\lambda}\delta^{\nu\rho})\Id - \epsilon^{\mu\nu\rho\lambda}\mathcal{D} + \mathcal{A}^{\mu\nu\rho\lambda},
\end{equation}
where $\mathcal{D} \equiv -c^0 c^1 c^2 c^3$ is the local fermion parity operator, $\epsilon^{\mu\nu\rho\lambda}$ is the totally antisymmetric Levi-Civita symbol ($\epsilon^{0123} = +1$), and $\mathcal{A}^{\mu\nu\rho\lambda}$ is an antisymmetric linear combination of quadratic Majorana bilinears:
\begin{equation}
\label{eq:m1_1_A_bilinears_explicit}
\mathcal{A}^{\mu\nu\rho\lambda} \equiv \delta^{\mu\nu} c^\rho c^\lambda - \delta^{\mu\rho} c^\nu c^\lambda + \delta^{\mu\lambda} c^\nu c^\rho + \delta^{\rho\lambda} c^\mu c^\nu - \delta^{\nu\lambda} c^\mu c^\rho + \delta^{\nu\rho} c^\mu c^\lambda.
\end{equation}
\end{theorem}
\begin{proof}
If all four indices $\mu, \nu, \rho, \lambda$ are mutually distinct, they form an even or odd permutation of $(0, 1, 2, 3)$. The Kronecker delta terms in Eq.~\eqref{eq:m1_1_four_majorana_contraction} and $\mathcal{A}^{\mu\nu\rho\lambda}$ vanish identically, leaving $c^\mu c^\nu c^\rho c^\lambda = \epsilon^{\mu\nu\rho\lambda} c^0 c^1 c^2 c^3 = -\epsilon^{\mu\nu\rho\lambda}\mathcal{D}$. If two or more indices coincide, repeatedly anticommuting $c^\mu c^\nu = 2\delta^{\mu\nu}\Id - c^\nu c^\mu$ reduces the product to pairs of quadratic bilinears and scalar identity terms matching Eq.~\eqref{eq:m1_1_A_bilinears_explicit}.
\end{proof}

\begin{remark}[Physical Utility of the Contraction Identity]
Theorem~\ref{thm:m1_1_four_majorana_contraction} is of paramount physical significance: it proves that \textbf{no local quartic Majorana interaction can be algebraically independent on a single site}. Any quartic on-site operator collapses identically into a scalar, a quadratic bilinear, and the local parity operator $\mathcal{D}$. Consequently, on-site many-body interactions cannot generate irreducible quartic correlations within a single site.
\end{remark}

% EDITORIAL QUESTION: How does the quadratic Majorana bilinear algebra realize the so(4) = su(2)_spin + su(2)_gauge Lie algebra homomorphism?
\section{Lie Algebra Homomorphism: \texorpdfstring{$\mathfrak{so}(4) \cong \mathfrak{su}(2)_{\mathrm{spin}} \oplus \mathfrak{su}(2)_{\mathrm{gauge}}$}{so(4)}}
\label{sec:m1_1_lie_homomorphism}
% CLAIM: K-003 (provenance: doc/rnd/phase_1/m1_1_t1_local_majorana_algebra_conventions.md, test: test/unit/test_local_majorana_algebra.py::test_t00a_05_gauge_spin_lie_algebras_and_mutual_commutation)

\begin{definition}[Majorana Quadratic Bilinear Map]
\label{def:m1_1_bilinear_map}
Let $\mathfrak{so}(4)$ be the Lie algebra of real $4 \times 4$ skew-symmetric matrices ($A^T = -A$). For any $A \in \mathfrak{so}(4)$, we define the Hermitian operator $\mathcal{O}_A$ by:
\begin{equation}
\label{eq:m1_1_bilinear_def}
\mathcal{O}_A \equiv \frac{\ii}{4} c^T A c = \frac{\ii}{4}\sum_{\mu, \nu=0}^3 A_{\mu\nu} c^\mu c^\nu.
\end{equation}
\end{definition}

\begin{theorem}[Lie Algebra Homomorphism]
\label{thm:m1_1_lie_homomorphism}
The linear mapping $A \mapsto \mathcal{O}_A$ from $\mathfrak{so}(4)$ into the space of Hermitian operators on $\mathcal{F}$ is a Lie algebra homomorphism:
\begin{equation}
\label{eq:m1_1_lie_homomorphism}
[\mathcal{O}_A, \mathcal{O}_B] = \ii \mathcal{O}_{[A, B]}, \qquad \forall A, B \in \mathfrak{so}(4),
\end{equation}
where $[A, B] = AB - BA$ is the matrix commutator in $\mathfrak{so}(4)$.
\end{theorem}
\begin{proof}
Directly compute the commutator $[\mathcal{O}_A, \mathcal{O}_B]$ using Definition~\ref{def:m1_1_bilinear_map}:
\begin{equation}
[\mathcal{O}_A, \mathcal{O}_B] = -\frac{1}{16}\sum_{\mu,\nu,\rho,\lambda=0}^3 A_{\mu\nu} B_{\rho\lambda} [c^\mu c^\nu, c^\rho c^\lambda].
\end{equation}
Using the CAR identity $[c^\mu c^\nu, c^\rho c^\lambda] = 2\delta^{\nu\rho} c^\mu c^\lambda - 2\delta^{\mu\rho} c^\nu c^\lambda + 2\delta^{\nu\lambda} c^\rho c^\mu - 2\delta^{\mu\lambda} c^\rho c^\nu$, we contract with the skew-symmetric tensors $A_{\mu\nu} = -A_{\nu\mu}$ and $B_{\rho\lambda} = -B_{\lambda\rho}$:
\begin{equation}
\sum_{\mu,\nu,\rho,\lambda} A_{\mu\nu} B_{\rho\lambda} [c^\mu c^\nu, c^\rho c^\lambda] = 4 \sum_{\mu,\lambda} (AB - BA)_{\mu\lambda} c^\mu c^\lambda = 4 c^T [A, B] c.
\end{equation}
Multiplying by $-\frac{1}{16}$ yields:
\begin{equation}
[\mathcal{O}_A, \mathcal{O}_B] = -\frac{1}{4} c^T [A, B] c = \ii \left( \frac{\ii}{4} c^T [A, B] c \right) = \ii \mathcal{O}_{[A, B]}.
\end{equation}
This completes the proof.
\end{proof}

The 6-dimensional Lie algebra $\mathfrak{so}(4)$ decomposes into the direct sum of two commuting 3-dimensional simple ideals: $\mathfrak{so}(4) \cong \mathfrak{su}(2)_{\mathrm{spin}} \oplus \mathfrak{su}(2)_{\mathrm{gauge}}$. The basis generators correspond to self-dual chiral matrices $\mathrm{M}^\gamma$ and anti-self-dual 't Hooft matrices $\mathrm{G}^\gamma$ ($\gamma \in \{1, 2, 3\}$):
\begin{align}
\label{eq:m1_1_G_matrices}
\mathrm{G}^1 &= \Id_2 \otimes \epsilon, & \mathrm{G}^2 &= \epsilon \otimes \sigma_z, & \mathrm{G}^3 &= \epsilon \otimes \sigma_x, \\
\label{eq:m1_1_M_matrices}
\mathrm{M}^1 &= \sigma_z \otimes \epsilon, & \mathrm{M}^2 &= \epsilon \otimes \Id_2, & \mathrm{M}^3 &= \sigma_x \otimes \epsilon,
\end{align}
where $\epsilon \equiv \ii \sigma_y = \begin{pmatrix} 0 & 1 \\ -1 & 0 \end{pmatrix}$. Their explicit contractions evaluate to:
\begin{align}
\label{eq:m1_1_G_contraction}
c^T \mathrm{G}^\gamma c &= 2(c^0 c^\gamma + c^{\gamma^+} c^{\gamma^-}), \\
\label{eq:m1_1_M_contraction}
c^T \mathrm{M}^\gamma c &= 2(c^0 c^\gamma - c^{\gamma^+} c^{\gamma^-}),
\end{align}
where $(\gamma, \gamma^+, \gamma^-)$ denotes cyclic permutations of $(1, 2, 3)$.

\begin{definition}[Spin and Gauge Generator Bilinears]
\label{def:m1_1_spin_gauge_bilinears}
The normalized physical spin and gauge bilinears on $\mathcal{F}$ are defined by:
\begin{equation}
\label{eq:m1_1_spin_gauge_bilinears}
S^\gamma \equiv \frac{1}{2} \mathcal{O}_{\mathrm{M}^\gamma} = \frac{\ii}{8} c^T \mathrm{M}^\gamma c, \qquad T^\gamma \equiv \frac{1}{2} \mathcal{O}_{\mathrm{G}^\gamma} = \frac{\ii}{8} c^T \mathrm{G}^\gamma c.
\end{equation}
\end{definition}

\begin{proposition}[Commutation Relations and Gauge Minus Sign]
\label{prop:m1_1_su2_algebras}
The generator matrices satisfy:
\begin{align}
\label{eq:m1_1_matrix_comm}
[\mathrm{M}^\alpha, \mathrm{M}^\beta] &= 2 \epsilon^{\alpha\beta\gamma} \mathrm{M}^\gamma, & [\mathrm{G}^\alpha, \mathrm{G}^\beta] &= -2 \epsilon^{\alpha\beta\gamma} \mathrm{G}^\gamma, & [\mathrm{M}^\alpha, \mathrm{G}^\beta] &= 0, \\
\label{eq:m1_1_trace_ortho}
\tr((\mathrm{M}^\alpha)^T \mathrm{M}^\beta) &= 4 \delta^{\alpha\beta}, & \tr((\mathrm{G}^\alpha)^T \mathrm{G}^\beta) &= 4 \delta^{\alpha\beta}, & \tr((\mathrm{M}^\alpha)^T \mathrm{G}^\beta) &= 0.
\end{align}
Consequently, by Theorem~\ref{thm:m1_1_lie_homomorphism}, the operator bilinears satisfy two mutually commuting $\mathfrak{su}(2)$ Lie algebras:
\begin{equation}
\label{eq:m1_1_op_comm}
[S^\alpha, S^\beta] = \ii \epsilon^{\alpha\beta\gamma} S^\gamma, \qquad [T^\alpha, T^\beta] = -\ii \epsilon^{\alpha\beta\gamma} T^\gamma, \qquad [S^\alpha, T^\beta] = 0.
\end{equation}
\end{proposition}
\begin{remark}
The minus sign in the gauge commutator $[T^\alpha, T^\beta] = -\ii\epsilon^{\alpha\beta\gamma} T^\gamma$ is an intrinsic feature of the anti-self-dual representation $\mathfrak{su}(2)_{\mathrm{gauge}} \subset \mathfrak{so}(4)$. Defining $J_{\mathrm{gauge}}^\alpha \equiv -T^\alpha$ restores the standard right-handed sign convention.
\end{remark}

% EDITORIAL QUESTION: How do the spin and gauge quadratic Casimirs relate to the local fermion parity operator, and how is the physical subspace projected?
\section{Casimir-Parity Invariants, Single-Generator Theorems, and Physical Projection}
\label{sec:m1_1_casimir_parity}
% CLAIM: K-002 (provenance: doc/rnd/phase_1/m1_1_t2_spin_gauge_representation.md, test: test/unit/test_local_majorana_representations.py::test_t00b_03_spin_and_gauge_casimirs_equal_parity_expressions)
The local fermion parity operator $\mathcal{D}$ is defined by:
\begin{equation}
\label{eq:m1_1_parity_def}
\mathcal{D} \equiv -c^0 c^1 c^2 c^3 = -(-1)^{n_1 + n_2}, \qquad \mathcal{D}^\dagger = \mathcal{D}, \quad \mathcal{D}^2 = \Id.
\end{equation}

\begin{theorem}[Single-Generator and Casimir Parity Identities]
\label{thm:m1_1_casimir_identities}
For each individual gauge generator $T^\gamma$ and spin generator $S^\gamma$ ($\gamma \in \{1, 2, 3\}$), the following single-generator operator identities hold exactly on $\mathcal{F}$:
\begin{equation}
\label{eq:m1_1_single_generator_parity}
(T^\gamma)^2 = \frac{1}{8}(\Id - \mathcal{D}), \qquad (S^\gamma)^2 = \frac{1}{8}(\Id + \mathcal{D}).
\end{equation}
Summing over $\gamma \in \{1, 2, 3\}$, the total quadratic Casimirs satisfy:
\begin{equation}
\label{eq:m1_1_casimir_sums}
\mathbf{T}^2 \equiv \sum_{\gamma=1}^3 (T^\gamma)^2 = \frac{3}{8}(\Id - \mathcal{D}), \qquad \mathbf{S}^2 \equiv \sum_{\gamma=1}^3 (S^\gamma)^2 = \frac{3}{8}(\Id + \mathcal{D}).
\end{equation}
\end{theorem}
\begin{proof}
Directly evaluate $(T^1)^2$ using $T^1 = \frac{\ii}{4}(c^0 c^1 + c^2 c^3)$:
\begin{equation}
(T^1)^2 = -\frac{1}{16}\left( (c^0 c^1)^2 + (c^2 c^3)^2 + c^0 c^1 c^2 c^3 + c^2 c^3 c^0 c^1 \right).
\end{equation}
Since $(c^0 c^1)^2 = c^0 c^1 c^0 c^1 = -(c^0)^2 (c^1)^2 = -\Id$, and $c^2 c^3 c^0 c^1 = c^0 c^1 c^2 c^3 = -\mathcal{D}$:
\begin{equation}
(T^1)^2 = -\frac{1}{16}(-\Id - \Id - 2\mathcal{D}) = \frac{1}{8}(\Id - \mathcal{D}).
\end{equation}
By cyclic symmetry of $(1, 2, 3)$, Eq.~\eqref{eq:m1_1_single_generator_parity} holds for all $\gamma$, and summing over $\gamma$ yields $\mathbf{T}^2 = \frac{3}{8}(\Id - \mathcal{D})$. For $S^1 = \frac{\ii}{4}(c^0 c^1 - c^2 c^3)$, the cross-term has the opposite sign, yielding $(S^1)^2 = \frac{1}{8}(\Id + \mathcal{D})$ and $\mathbf{S}^2 = \frac{3}{8}(\Id + \mathcal{D})$.
\end{proof}

The 4-dimensional local Fock space $\mathcal{F}$ decomposes into two orthogonal parity eigenspaces:
\begin{equation}
\label{eq:m1_1_fock_decomposition}
\mathcal{F} = \mathcal{F}_+ \oplus \mathcal{F}_- \cong \left(\tfrac{1}{2}, 0\right)_{\mathrm{spin,gauge}} \oplus \left(0, \tfrac{1}{2}\right)_{\mathrm{spin,gauge}}.
\end{equation}
The physical spin-$1/2$ Hilbert space is the gauge-singlet sector $\mathcal{F}_+ = \ker(\mathcal{D} - \Id) = \bigcap_\gamma \ker T^\gamma$, on which $\mathbf{S}^2 = \frac{3}{4}\Id$ and $\mathbf{T}^2 = 0$. The exact Gutzwiller physical projector is:
\begin{equation}
\label{eq:m1_1_gutzwiller_projector}
P_G = \frac{\Id + \mathcal{D}}{2}.
\end{equation}

% EDITORIAL QUESTION: How does physical projection eliminate gauge leakage to recover the standard SU(2) Pauli algebra?
\section{Physical Spin Operators and the Pauli Algebra Modulo Gauge}
\label{sec:m1_1_pauli_algebra}
In the Kitaev Majorana representation, physical spin operators are represented through:
\begin{equation}
\label{eq:m1_1_pauli_def}
\sigma^\gamma \equiv -\ii c^0 c^\gamma = \frac{\ii}{2} c^T \mathrm{N}^\gamma c, \qquad \mathrm{N}^\gamma \equiv \frac{1}{2}(\mathrm{M}^\gamma + \mathrm{G}^\gamma),
\end{equation}
where $\mathrm{N}^1 = e_{11} \otimes \epsilon$, $\mathrm{N}^2 = \epsilon \otimes e_{11}$, and $\mathrm{N}^3 = \frac{1}{2}(\epsilon \otimes \sigma_x + \sigma_x \otimes \epsilon)$, with $e_{11} = \frac{1}{2}(\Id_2 + \sigma_z)$.

\begin{proposition}[Pauli Algebra Modulo Gauge]
\label{prop:m1_1_pauli_modulo_gauge}
For cyclic indices $(\alpha, \beta, \gamma)$:
\begin{equation}
\label{eq:m1_1_pauli_cyclic}
\sigma^\alpha \sigma^\beta = \ii \sigma^\gamma - 4\ii T^\gamma, \qquad \sigma^\gamma = 2(S^\gamma + T^\gamma).
\end{equation}
For general $\alpha, \beta \in \{1, 2, 3\}$, the operator product expands as:
\begin{equation}
\label{eq:m1_1_pauli_general}
\sigma^\alpha \sigma^\beta = \delta^{\alpha\beta}\Id + \ii \sum_{\gamma=1}^3 \epsilon^{\alpha\beta\gamma}(\sigma^\gamma - 4 T^\gamma).
\end{equation}
\end{proposition}
\begin{proof}
For $\alpha \neq \beta$, $\sigma^\alpha \sigma^\beta = (-\ii c^0 c^\alpha)(-\ii c^0 c^\beta) = c^\alpha c^\beta$. Using Eq.~\eqref{eq:m1_1_G_contraction}, $c^\alpha c^\beta = \frac{1}{2}c^T \mathrm{G}^\gamma c - c^0 c^\gamma = -4\ii T^\gamma + \ii \sigma^\gamma$. For $\alpha = \beta$, $(\sigma^\alpha)^2 = -(c^0 c^\alpha)^2 = \Id$.
\end{proof}

Sandwiching Eq.~\eqref{eq:m1_1_pauli_general} with the Gutzwiller projector $P_G$ and utilizing $T^\gamma P_G = 0$:
\begin{equation}
\label{eq:m1_1_projected_pauli}
P_G \sigma^\alpha \sigma^\beta P_G = \delta^{\alpha\beta} P_G + \ii \sum_{\gamma=1}^3 \epsilon^{\alpha\beta\gamma} P_G \sigma^\gamma P_G.
\end{equation}
Thus, on the physical subspace $\Ran P_G$, the operators $\sigma_{\mathrm{phys}}^\gamma \equiv P_G \sigma^\gamma P_G$ satisfy the exact $\mathrm{SU}(2)$ Pauli spin algebra with zero gauge leakage.

% EDITORIAL QUESTION: How is the microscopic interaction formulated on a single bond before assembling the full lattice Hamiltonian?
\section{Single-Bond Interaction and the 5-Parameter Exchange Tensor}
\label{sec:m1_1_single_bond}
Before treating the general many-body lattice problem, we analyze the microscopic exchange interaction on an isolated two-site bond connecting site $i$ (sublattice $A$) and site $j$ (sublattice $B$).

Let the bond carry spatial orientation $\gamma \in \{x, y, z\}$. In the presence of strong spin-orbit coupling, the spin exchange interaction is described by the 5-parameter exchange tensor $\mathbf{J}_\gamma$:
\begin{equation}
\label{eq:m1_1_single_bond_H}
H_{ij}^{(\gamma)} = \sum_{a, b = 1}^3 \sigma_i^a \mathbf{J}_\gamma^{ab} \sigma_j^b.
\end{equation}

\begin{definition}[Extended 5-Parameter Exchange Tensor]
\label{def:m1_1_exchange_tensor}
On a bond of type $\gamma$ with cyclic labels $(\gamma, \gamma^+, \gamma^-)$, the exchange tensor $\mathbf{J}_\gamma$ decomposes into symmetric Heisenberg ($J$), Kitaev ($K$), symmetric off-diagonal ($\Gamma, \Gamma'$), and antisymmetric Dzyaloshinskii-Moriya ($D$) interactions:
\begin{equation}
\label{eq:m1_1_J_tensor_components}
\begin{aligned}
\mathbf{J}_\gamma^{ab} ={}& J \delta^{ab} + K \delta^{a\gamma} \delta^{b\gamma} + \Gamma (\delta^{a\gamma^+} \delta^{b\gamma^-} + \delta^{a\gamma^-} \delta^{b\gamma^+}) \\
&+ \Gamma' (\delta^{a\gamma^+} \delta^{b\gamma} + \delta^{a\gamma} \delta^{b\gamma^+} + \delta^{a\gamma^-} \delta^{b\gamma} + \delta^{a\gamma} \delta^{b\gamma^-}) \\
&+ D (\delta^{a\gamma^+} \delta^{b\gamma^-} - \delta^{a\gamma^-} \delta^{b\gamma^+}).
\end{aligned}
\end{equation}
In the canonical Cartesian basis $(x, y, z)$, the $z$-bond tensor ($\gamma = 3$) has the explicit matrix form:
\begin{equation}
\label{eq:m1_1_Jz_matrix}
\mathbf{J}_z = \begin{pmatrix} J & \Gamma + D & \Gamma' \\ \Gamma - D & J & \Gamma' \\ \Gamma' & \Gamma' & J + K \end{pmatrix}.
\end{equation}
\end{definition}

\begin{remark}[Bond Transposition and Edge Reversal]
Under edge reversal $j \to i$, the interaction matrix is transposed: $\mathbf{J}_{ji} = \mathbf{J}_{ij}^T$. The symmetric couplings $J, K, \Gamma, \Gamma'$ remain invariant, whereas the DM coupling flips sign: $D \to -D$. This guarantees exact angular momentum conservation across directed edges.
\end{remark}

Substituting the Kitaev representation $\sigma_i^a = \frac{\ii}{2} c_i^T \mathrm{N}^a c_i$ and $\sigma_j^b = \frac{\ii}{2} c_j^T \mathrm{N}^b c_j$ into Eq.~\eqref{eq:m1_1_single_bond_H}, the single-bond interaction becomes an exact quartic Majorana operator:
\begin{equation}
\label{eq:m1_1_bond_quartic}
H_{ij}^{(\gamma)} = -\frac{1}{4} \sum_{a,b=1}^3 \mathbf{J}_\gamma^{ab} (c_i^T \mathrm{N}^a c_i)(c_j^T \mathrm{N}^b c_j).
\end{equation}

% EDITORIAL QUESTION: How is the full lattice Hamiltonian assembled, and why is the projected operator equivalence exact?
\section{Lattice Assembly, Gauge Multipliers, and Projected Equivalence}
\label{sec:m1_1_quartic_hamiltonian}
We now assemble the full lattice Hamiltonian on a bipartite honeycomb graph $\Lambda = (\mathcal{V}, \mathcal{E})$. We partition the Hamiltonian into vertex terms (Zeeman fields and local gauge constraints) and directed bond terms:
\begin{equation}
\label{eq:m1_1_spin_hamiltonian}
H = \sum_{b=(i,j) \in \overrightarrow{\mathcal{E}}_{A \to B}} H_{ij}^{(\gamma_b)} - \sum_{i \in \mathcal{V}} \sum_{a=1}^3 h_i^a \sigma_i^a.
\end{equation}
Introducing local Lagrange multiplier terms $\lambda_i^\gamma T_i^\gamma$ at each vertex $i \in \mathcal{V}$ to enforce gauge-singlet expectation values $\langle T_i^\gamma \rangle = 0$, the extended Majorana Hamiltonian reads:
\begin{equation}
\label{eq:m1_1_H_lambda_full}
\begin{aligned}
H_\lambda ={}& \frac{1}{8} \sum_{(i,j) \in \overrightarrow{\mathcal{E}}} \sum_{a,b=1}^3 (\ii c_i^T \mathrm{N}^a c_i) \mathbf{J}_{\gamma}^{ab} (\ii c_j^T \mathrm{N}^b c_j) \\
&- \frac{1}{2} \sum_{i \in \mathcal{V}} \sum_{a=1}^3 h_i^a (\ii c_i^T \mathrm{N}^a c_i) - \sum_{i \in \mathcal{V}} \sum_{\gamma=1}^3 \lambda_i^\gamma T_i^\gamma.
\end{aligned}
\end{equation}
Here, the lattice Gutzwiller projector is the tensor product of local projectors:
\begin{equation}
\label{eq:m1_1_global_gutzwiller}
P_G = \prod_{i \in \mathcal{V}} P_{G, i} = \prod_{i \in \mathcal{V}} \frac{\Id + \mathcal{D}_i}{2}.
\end{equation}

\begin{theorem}[Exact Projected Operator Equivalence]
\label{thm:m1_1_projected_equivalence}
Let $P_G$ be the global Gutzwiller projector onto the physical parity-even gauge-singlet sector. Then $H_\lambda$ and the microscopic spin Hamiltonian $H$ satisfy the exact operator equivalence:
\begin{equation}
\label{eq:m1_1_projected_equivalence}
P_G H_\lambda P_G = P_G H P_G.
\end{equation}
\end{theorem}
\begin{proof}
On the physical subspace, $T_i^\gamma P_G = 0$ for all $i \in \mathcal{V}$ and $\gamma \in \{1, 2, 3\}$. Hence, the Lagrange multiplier term vanishes identically: $\left(\sum_{i,\gamma}\lambda_i^\gamma T_i^\gamma\right) P_G = 0$. Furthermore, by Proposition~\ref{prop:m1_1_pauli_modulo_gauge}, $P_G \sigma_i^a P_G = P_G (\frac{\ii}{2} c_i^T \mathrm{N}^a c_i) P_G$. Since distinct site operators commute, the product of Majorana bilinears projected with $P_G$ evaluates identically to the physical spin exchange term.
\end{proof}

\section{Summary and Theoretical Transition}
\label{sec:m1_1_summary}
The fractionalization of spin-$1/2$ degrees of freedom into four Majorana fermions provides an exact microscopic representation of the extended Kitaev Hamiltonian via Theorem~\ref{thm:m1_1_projected_equivalence}. The representation is rigorous on the operator level.

However, the Hamiltonian $H_\lambda$ is **quartic** in Majorana operators. Solving this interacting fermion problem on the full many-body Hilbert space is intractable. In the subsequent chapter (\Cref{ch:m1_1b}), we develop the model-independent theory of **quasi-free states**, showing that mean-field decoupling is not an operator approximation, but the exact restriction of the quantum variational principle to the fermionic Grassmannian manifold.
